\epanalipsi{Β4}

\tmima{Λεξικό του κεφαλαίου}

\renewcommand{\arraystretch}{1.5}
\begin{longtable}{@{}>{\raggedright\arraybackslash}p{0.28\linewidth} >{\raggedright\arraybackslash}p{0.62\linewidth}@{}}
\toprule
\textbf{Λέξη ή σύμβολο} & \textbf{Με μία πρόταση}\\
\midrule
\endhead
έδρα, ακμή, κορυφή & Επίπεδη επιφάνεια του στερεού· τομή δύο εδρών· σημείο συνάντησης ακμών.\\
ορθό πρίσμα & Δύο ίσες παράλληλες βάσεις (πολύγωνα) και παράπλευρες έδρες ορθογώνια.\\
παραλληλεπίπεδο, κύβος & Πρίσμα με βάσεις παραλληλόγραμμα· ορθογώνιο παραλληλεπίπεδο με όλες τις έδρες τετράγωνα.\\
πυραμίδα & Μία βάση (πολύγωνο) και παράπλευρες έδρες τρίγωνα με κοινή κορυφή.\\
κύλινδρος & Παράγεται από ορθογώνιο που περιστρέφεται γύρω από μία πλευρά του· δύο κυκλικές βάσεις.\\
κώνος & Παράγεται από ορθογώνιο τρίγωνο που περιστρέφεται γύρω από μία κάθετη πλευρά του· μία κυκλική βάση.\\
σφαίρα & Παράγεται από κύκλο που περιστρέφεται γύρω από μία διάμετρό του.\\
ανάπτυγμα & Το στερεό ξεδιπλωμένο στο επίπεδο.\\
όψη & Το σχέδιο του στερεού από μπροστά, από πάνω ή από το πλάι.\\
$\nu$-γωνικό πρίσμα & $\nu+2$ έδρες, $2\nu$ κορυφές, $3\nu$ ακμές.\\
$\nu$-γωνική πυραμίδα & $\nu+1$ έδρες, $\nu+1$ κορυφές, $2\nu$ ακμές.\\
\bottomrule
\end{longtable}

\tmima{Ασκήσεις επανάληψης}

\ex Ονόμασε το στερεό που θυμίζει καθένα: α) κουτί γάλακτος \quad β) τενεκεδάκι με φασόλια \quad γ) κορυφή σκεπής σε σχήμα πυραμίδας \quad δ) φρυγανιά σε σχήμα κουτιού \quad ε) χωνάκι παγωτού \quad στ) μήλο σε σχήμα σφαίρας.

\ex Βρες τις έδρες, τις ακμές και τις κορυφές: α) του κύβου \quad β) του τετραέδρου \quad γ) του τριγωνικού πρίσματος.

\ex Ένα πρίσμα έχει βάση οκτάγωνο. Πόσες έδρες, ακμές και κορυφές έχει;

\ex Μια πυραμίδα έχει $20$ ακμές. Τι βάση έχει και πόσες έδρες και κορυφές;

\ex Ποια από τις διατάξεις έξι τετραγώνων είναι ανάπτυγμα κύβου; α) ένας «σταυρός» με στήλη τεσσάρων τετραγώνων και ένα τετράγωνο αριστερά και δεξιά του δεύτερου \quad β) δύο σειρές των τριών τετραγώνων, η μία πάνω από την άλλη \quad γ) έξι τετράγωνα σε μία σειρά.

\ex Αντιστοίχισε κάθε ανάπτυγμα με το στερεό του: α) ορθογώνιο και δύο κύκλοι \quad β) κυκλικός τομέας και κύκλος \quad γ) τετράγωνο και τέσσερα ίσα ισοσκελή τρίγωνα \quad δ) τρία ορθογώνια σε σειρά και δύο ισόπλευρα τρίγωνα.

\ex Ένα κουτί σε σχήμα ορθογωνίου παραλληλεπιπέδου έχει διαστάσεις $5$, $3$ και $2$ cm. Πόσο είναι το άθροισμα των μηκών όλων των ακμών του;

\ex Κατασκεύασε το ανάπτυγμα κυλίνδρου με ακτίνα βάσης $1{,}5$ cm και ύψος $4$ cm. Πόσο είναι το μήκος του ορθογωνίου; (Πάρε $\pi\approx3{,}14$.)

\ex Ένας κώνος έχει ακτίνα βάσης $3$ cm και γενέτειρα $6$ cm. Βρες τη γωνία του κυκλικού τομέα του αναπτύγματός του. Τι παρατηρείς για τον τομέα;

\ex Ποιο στερεό παράγεται όταν περιστρέψουμε: α) ορθογώνιο γύρω από μία πλευρά του \quad β) ορθογώνιο τρίγωνο γύρω από μία κάθετη πλευρά του \quad γ) κύκλο γύρω από μια διάμετρό του;

\ex Σχεδίασε την όψη από μπροστά και την όψη από πάνω ενός κυλίνδρου που στέκεται όρθιος.

\ex Η σφαίρα έχει διάμετρο $18$ cm. Πόση είναι η ακτίνα της; Πόσες έδρες έχει;

\ex Ποιο στερεό έχει $8$ έδρες και $12$ κορυφές; Πόσες ακμές έχει;

\ex Από πάνω ένα στερεό φαίνεται σαν τετράγωνο με τις δύο διαγωνίους του και από μπροστά σαν τρίγωνο. Ποιο στερεό είναι;

\ex \textbf{Σωστό ή λάθος (πολύεδρα).} α) Ο κύβος είναι ορθογώνιο παραλληλεπίπεδο. \quad β) Η τετραγωνική πυραμίδα έχει $5$ κορυφές. \quad γ) Κάθε πρίσμα έχει παράπλευρες έδρες τρίγωνα. \quad δ) Το τετράεδρο έχει $4$ έδρες και $6$ ακμές.

\ex \textbf{Σωστό ή λάθος (καμπύλα στερεά).} α) Ο κώνος έχει δύο βάσεις. \quad β) Η σφαίρα δεν έχει ανάπτυγμα. \quad γ) Η γενέτειρα του κώνου είναι μεγαλύτερη από το ύψος του. \quad δ) Το ανάπτυγμα του κυλίνδρου έχει ένα ορθογώνιο και δύο κύκλους.

\Needspace{17\baselineskip}\tmima{Ο χάρτης της επανάληψης}

Μετά τις λύσεις, δες σε ποιες ασκήσεις δυσκολεύτηκες. Ο πίνακας σου λέει ποιο μέρος της παραγράφου να ξαναδιαβάσεις.

\renewcommand{\arraystretch}{1.5}
\begin{tabular*}{\linewidth}{@{\extracolsep{\fill}} >{\raggedright\arraybackslash}p{0.38\linewidth} >{\raggedright\arraybackslash}p{0.52\linewidth}}
\toprule
\textbf{Αν δυσκολεύτηκες στις} & \textbf{ξαναδιάβασε}\\
\midrule
ΕΒ4.1, ΕΒ4.2, ΕΒ4.3, ΕΒ4.7, ΕΒ4.15 & 4.1 Α (έδρες, ακμές, κορυφές, πρίσμα, κύβος)\\
ΕΒ4.4, ΕΒ4.13 & 4.1 Β (πυραμίδα)\\
ΕΒ4.5, ΕΒ4.6 & 4.1 Α και Β (αναπτύγματα)\\
ΕΒ4.8, ΕΒ4.9, ΕΒ4.10, ΕΒ4.12, ΕΒ4.16 & 4.1 Γ (κύλινδρος, κώνος, σφαίρα)\\
ΕΒ4.11, ΕΒ4.14 & 4.1 Δ (όψεις)\\
\bottomrule
\end{tabular*}
